\section{Domain Task 8: Phân tích mùa vụ và cơ hội đặt phòng}

Câu hỏi: Tỷ lệ lấp đầy thay đổi như thế nào theo mùa? Chính sách minimum
nights ảnh hưởng thế nào?

\subsection{Biểu đồ 1 -- Line Chart: Tỷ lệ lấp đầy theo tháng}

\subsubsection*{A. Thiết kế Idiom}

\begin{longtable}{|p{2.8cm}|p{11.5cm}|}
\hline
\textbf{Đặc điểm} & \textbf{Chi tiết} \\
\hline
\endhead
\hline
\endfoot
Idiom & Line Chart \\
\hline
What & Month: Ordinal (tháng 1--12, cyclic -- seasonality lặp lại theo năm) \newline
occupancy\_rate\_pct (derived): Quantitative (AVG(is\_booked) $\times$ 100) \\
\hline
How & \textbf{Encode:} \newline
\newline
Mark: Line + Point \newline
Channel: \newline
\hspace{0.3cm} Pos X: Month (Ordinal, Discrete -- tiến trình thời gian tháng 1--12) \newline
\hspace{0.3cm} Pos Y: Occupancy Rate (\%) (Quantitative) \newline
\hspace{0.3cm} Reference Line: Average (ngang -- mức trung bình tổng thể $\sim$30\%) \newline
\newline
\textbf{Manipulate:} \newline
\newline
Hover để xem tháng, occupancy rate (\%) cụ thể \newline
\newline
\textbf{Facet:} \newline
\newline
Superimpose: Reference Line (average occupancy ngang) \newline
\newline
\textbf{Reduce:} \newline
\newline
Filter: Year = 2026 (dữ liệu calendar 2025 chỉ có từ tháng 11, không đủ để phân tích seasonality) \\
\hline
Why & discover $\rightarrow$ browse \newline
Line chart (không phải bar) vì Month là Ordinal có thứ tự -- đường kết nối nhấn mạnh tính liên tục và xu hướng thời gian, không phải giá trị tại từng điểm rời rạc. \\
\hline
Scale & Key: 12 (tháng); Items: dữ liệu 2026 (tháng 1--11) \\
\hline
\end{longtable}

\begin{figure}[H]
    \centering
    \safeincludegraphics[width=0.9\linewidth]{images/domain_task8_chart1.png}
    \caption{Hình 8.1. Tỷ lệ lấp đầy (occupancy rate) theo tháng tại NYC (2026)}
    \label{fig:domain-task8-chart1}
\end{figure}

\subsubsection*{B. Đánh giá biểu đồ}

\textbf{1. Nguyên lý biểu đạt (Expressiveness):}

\begin{itemize}
    \item Thuộc tính: Tháng / Month (O), Tỷ lệ lấp đầy / Occupancy Rate Pct (Q -- AVG aggregate), Đường trung bình / Average Occupancy Rate (Q -- Reference Line), Năm / Year (O -- filter bối cảnh, Year = 2026).
    \item Channel:
    \begin{itemize}
        \item PosX (Month -- Ordinal): trục thời gian có thứ tự $\rightarrow$ đúng. Không dùng Categorical vì mất thứ tự thời gian.
        \item PosY (Occupancy Rate -- Q): thể hiện mức độ lấp đầy $\rightarrow$ đúng.
        \item Reference Line Average: cung cấp mốc tham chiếu tổng thể $\rightarrow$ tăng thêm thông tin phân tích.
    \end{itemize}
    \item Đánh giá mức độ quan trọng:
    \begin{itemize}
        \item Ưu tiên 1 -- Vị trí dọc (PosY --- Position on common scale): PosY là kênh quan trọng nhất trong line chart. Kênh này mã hóa Occupancy Rate (\%) (Quantitative) trên thang đo chung, phản ánh trực tiếp mức độ lấp đầy của thị trường tại từng thời điểm. Theo Mackinlay, Position on a common scale là kênh hiệu quả nhất cho dữ liệu định lượng, cho phép người xem so sánh mức occupancy giữa các tháng với độ chính xác cao và nhận diện biên độ dao động mùa vụ (từ $\sim$20\% đến $\sim$40\%) một cách tức thì. Kênh này là nền tảng để thực hiện task discover (phát hiện xu hướng) và browse (duyệt qua từng tháng).
        \item Ưu tiên 2 -- Vị trí ngang (PosX --- trục thời gian Ordinal): PosX mã hóa Month (Ordinal --- tháng 1 đến 12) theo trục thời gian. Đây là kênh quan trọng thứ hai vì nó thiết lập chiều thời gian của biểu đồ --- nền tảng của phân tích seasonality. Theo Mackinlay, Position phù hợp cho dữ liệu thứ tự (Ordinal), và việc sử dụng trục ngang liên tục phản ánh đúng bản chất thứ tự tháng (T1 $\rightarrow$ T12), không bị mất thông tin thứ tự như khi dùng Categorical. Các điểm được kết nối bằng đường thẳng giúp người xem nhận diện xu hướng liên tục theo thời gian.
        \item Ưu tiên 3 -- Reference Line (Average): Reference line ngang tại mức trung bình ($\sim$30\%) là channel bổ trợ giúp người xem định vị nhanh các tháng vượt/dưới ngưỡng bình thường mà không cần đọc chính xác giá trị trục Y. Đây là annotation tĩnh (không encode data mới) nhưng tăng đáng kể khả năng thực hiện action Locate Extremes --- xác định tháng cao điểm và thấp điểm một cách tức thì. Theo phân loại Munzner, đây là dạng derived encoding bổ trợ cho channel chính thay vì thay thế nó.
    \end{itemize}
    \item Kết luận: Line chart là idiom tối ưu cho dữ liệu chuỗi thời gian, đặc biệt khi cần phát hiện trend và seasonality.
\end{itemize}

\textbf{2. Nguyên lý hiệu quả (Effectiveness):}

\begin{itemize}
    \item Độ chính xác (Accuracy): PosY (position on common scale) --- rất chính xác. Người xem dễ dàng so sánh mức occupancy giữa các tháng.
    \item Discriminability: Reference line Average tạo mốc tham chiếu để xác định tháng nào trên/dưới trung bình.
    \item Separability: PosX và PosY tách biệt hoàn toàn.
\end{itemize}

\subsubsection*{C. Phân tích biểu đồ (Insight)}

\begin{itemize}
    \item Tháng 1--4 là low season rõ rệt (occupancy $\sim$20--21\%) --- cơ hội tốt cho khách muốn giá thấp và dễ tìm phòng.
    \item Từ tháng 5 trở đi, occupancy tăng dần và đạt đỉnh vào tháng 11 ($\sim$40\%) --- phản ánh nhu cầu du lịch mùa thu và dịp lễ Thanksgiving.
    \item Dữ liệu sử dụng: calendar 2026 (tháng 1--11); dữ liệu 2025 chỉ có từ tháng 11 nên không đủ để so sánh liên năm.
    \item Khuyến nghị cho du khách: Đặt phòng tháng 1--3 để có giá thấp nhất; tháng 11 cần đặt sớm do cầu rất cao.
    \item Khuyến nghị cho host: Tháng 1--4 nên linh hoạt hóa chính sách (giảm minimum nights, giảm giá nhẹ); tháng 11--12 có thể tăng giá.
\end{itemize}

\subsection{Biểu đồ 2 -- Heatmap: Tỷ lệ lấp đầy theo tháng và chính sách minimum nights}

\subsubsection*{A. Thiết kế Idiom}

\begin{longtable}{|p{2.8cm}|p{11.5cm}|}
\hline
\textbf{Đặc điểm} & \textbf{Chi tiết} \\
\hline
\endhead
\hline
\endfoot
Idiom & Heatmap (Matrix Chart) \\
\hline
What & Month: Ordinal (tháng 1--12, cyclic) \newline
minimum\_nights\_group (derived): Categorical (Short $\leq$3 / Medium 4--7 / Long $>$7) \newline
occupancy\_rate\_pct (derived): Quantitative \\
\hline
How & \textbf{Encode:} \newline
\newline
Mark: Point (Được hiển thị dưới dạng vùng diện tích hình vuông -- Square/Area lấp đầy ô lưới ma trận) \newline
Channel: \newline
\hspace{0.3cm} Pos X: Month (Ordinal -- 12 cột) \newline
\hspace{0.3cm} Pos Y: Minimum Nights Group (Categorical -- Short/Medium/Long, top$\rightarrow$bottom) \newline
\hspace{0.3cm} Color (Luminance/Sequential): Occupancy Rate (Quantitative -- đậm=cao, nhạt=thấp) \newline
\newline
\textbf{Manipulate:} \newline
\newline
Hover để xem month, Minimum Nights Group, occupancy\_rate\_pct (\%) \newline
\newline
\textbf{Reduce:} \newline
\newline
Binning minimum\_nights $\rightarrow$ 3 nhóm (giảm số chiều từ hàng trăm giá trị $\rightarrow$ 3 nhóm) \newline
Sort thủ công: Short $\rightarrow$ Medium $\rightarrow$ Long (thứ tự logic nghiệp vụ, không alphabetical) \\
\hline
Why & discover $\rightarrow$ compare \newline
Phát hiện pattern 2 chiều đồng thời (month $\times$ minimum\_nights\_group). Heatmap theo nguyên lý `2 keys $\rightarrow$ heatmap' --- encode 3 biến trong cấu trúc ma trận. Color Luminance/Sequential đúng cho quantitative trong heatmap vì kích thước ô đồng đều loại bỏ bias từ Area. \\
\hline
Scale & Key: 12 (tháng) $\times$ 3 (nhóm) = 36 ô \newline
Bin: 3 (Short $\leq$3, Medium 4--7, Long $>$7 đêm) \\
\hline
\end{longtable}

\begin{figure}[H]
    \centering
    \safeincludegraphics[width=0.9\linewidth]{images/domain_task8_chart2.png}
    \caption{Hình 8.2. Heatmap tỷ lệ lấp đầy theo tháng và nhóm minimum nights}
    \label{fig:domain-task8-chart2}
\end{figure}

\subsubsection*{B. Đánh giá biểu đồ}

\textbf{1. Nguyên lý biểu đạt (Expressiveness):}

\begin{itemize}
    \item Thuộc tính: Tháng / Month (O), Nhóm đêm tối thiểu / Minimum Nights Group (C -- derived: Short $\leq$3 / Medium 4--7 / Long $>$7), Tỷ lệ lấp đầy / Occupancy Rate Pct (Q -- Color Luminance).
    \item Channel:
    \begin{itemize}
        \item PosX (Month -- O): trục thời gian theo thứ tự $\rightarrow$ đúng.
        \item PosY (Minimum Nights Group -- C): phân loại chính sách $\rightarrow$ đúng. Spatial region cho categorical.
        \item Color (Luminance/Sequential): thể hiện độ lớn occupancy (Q) $\rightarrow$ phù hợp. Sequential colormap (nhạt$\rightarrow$đậm) đúng cho quantitative attribute có hướng sequential.
    \end{itemize}
    \item Đánh giá mức độ quan trọng:
    \begin{itemize}
        \item Ưu tiên 1 -- Vị trí ngang (PosX --- chiều thời gian Ordinal): PosX mã hóa Month (Ordinal --- tháng 1 đến 12) dọc theo trục ngang của ma trận heatmap. Đây là kênh quan trọng nhất vì chiều thời gian là trục phân tích chính của task (seasonality). Theo Mackinlay, Position là kênh mạnh nhất cho dữ liệu thứ tự (Ordinal), đảm bảo thứ tự tháng được duy trì chính xác và người xem có thể theo dõi sự thay đổi của occupancy theo thời gian một cách tự nhiên từ trái sang phải --- phù hợp với convention đọc phương Tây và nhận thức thời gian tuyến tính.
        \item Ưu tiên 2 -- Vị trí dọc (PosY --- Spatial Region cho Categorical): PosY mã hóa Minimum Nights Group (Categorical --- Short/Medium/Long stay) theo hàng của ma trận. Đây là kênh quan trọng thứ hai vì nó xác định chiều so sánh thứ hai của biểu đồ: người xem so sánh giữa các chính sách minimum nights theo trục dọc. Theo Mackinlay, Spatial Region (phân vùng không gian) phù hợp cho dữ liệu danh mục --- mỗi hàng heatmap đóng vai trò như một spatial region riêng biệt, phân tách rõ ràng ba nhóm chính sách mà không cần dùng màu sắc (vốn đã được dùng cho kênh thứ ba).
        \item Ưu tiên 3 -- Độ sáng màu sắc (Color Luminance / Sequential): Color Luminance mã hóa Occupancy Rate (Quantitative) thông qua thang màu tuần tự (nhạt $\rightarrow$ đậm). Theo Mackinlay, Color Luminance không phải kênh tối ưu nhất cho dữ liệu định lượng về độ chính xác tuyệt đối (kém hơn Position hay Length), nhưng trong heatmap --- nơi cả PosX và PosY đã được dùng để mã hóa hai biến phân loại/thứ tự --- đây là kênh duy nhất khả dụng để mã hóa giá trị định lượng thứ ba trong không gian hai chiều. Color Luminance phù hợp để phát hiện pattern (nhận biết ô đậm/nhạt) và so sánh tương đối giữa các ô, phục vụ đúng mục tiêu discover của heatmap. Kích thước ô đồng đều (square) loại bỏ nhiễu từ kênh Area, giúp người xem chỉ tập trung vào màu sắc.
    \end{itemize}
    \item Kết luận: Heatmap là idiom phù hợp nhất khi cần phân tích 3 biến (2 keys + 1 value) trong cấu trúc ma trận.
\end{itemize}

\textbf{2. Nguyên lý hiệu quả (Effectiveness):}

\begin{itemize}
    \item Độ chính xác (Accuracy): Color (Luminance) --- accuracy thấp hơn Position, khó ước lượng chính xác \% chênh lệch. Tuy nhiên mục tiêu là phát hiện pattern (cao/thấp), không phải đọc giá trị chính xác --- trade-off chấp nhận được.
    \item Khả năng phân biệt (Discriminability): 36 ô vuông đều nhau. Sequential palette giúp phân biệt $\sim$5--7 cấp độ màu.
    \item Khả năng tách biệt (Separability): PosX, PosY và Color kết hợp tốt. Area đồng đều không tạo bias về kích thước.
\end{itemize}

\subsubsection*{C. Phân tích biểu đồ (Insight)}

\begin{itemize}
    \item Short minimum stay ($\leq$3 đêm) có occupancy cao và ổn định quanh năm --- chính sách linh hoạt thu hút nhiều khách nhất, đặc biệt trong mùa thấp điểm.
    \item Long minimum stay ($>$7 đêm) có occupancy thấp hơn trong phần lớn các tháng, nhưng tăng vào cuối năm --- phù hợp cho khách công tác dài hạn.
    \item Tháng 11--12 có màu đậm nhất ở hầu hết các nhóm --- xác nhận high season với nhu cầu cao bất kể chính sách minimum nights.
    \item Khuyến nghị cho host: Tháng 1--4 nên chuyển sang Short minimum stay; tháng 11--12 có thể giữ Medium/Long stay vì cầu cao tự nhiên bù đắp tính hạn chế của chính sách.
\end{itemize}
